\section{Front-end e Controle de Acesso}
\label{sec:front-end-controle-acesso}

As interfaces do módulo foram construídas sobre a arquitetura existente do Banco de
Especialistas, baseada em Vue.js, Nuxt e Vuetify. A seção apresenta os fundamentos dessa
camada --- reatividade e composição do Vue.js, organização de aplicações universais no
Nuxt, gerenciamento de estado com Vuex e controle de acesso baseado em papéis (RBAC)
---, pois decisões como a separação entre listagem administrativa e listagem do
solicitante, o uso de módulos Vuex por domínio e a exibição condicional do painel de
revisão decorrem dessas tecnologias e do modelo de segurança adotado.

\subsection{Vue.js}
\label{subsec:vue-js}

O Vue.js é um \textit{framework} progressivo, adotável de forma incremental
\cite{vuejs2024}, o que permitiu incorporar o módulo ao front-end legado sem reescrever
a camada de apresentação, em coerência com o \textit{Additive Change}
(Seção~\ref{subsec:additive-change}). Seu núcleo é o sistema de reatividade: o estado
declarado no componente é vinculado à interface e, quando alterado, o Vue atualiza
apenas as partes dependentes por meio de um \textit{virtual DOM}. No módulo, isso
permite refletir o retorno de uma decisão de revisão sem recarregar a página
(Seção~\ref{subsubsec:frontend-detalhe}).

A interface é organizada em componentes, sobretudo \textit{Single-File Components}
(\texttt{.vue}), com \texttt{<template>}, \texttt{<script>} e \texttt{<style>}. Essa
coesão permite isolar o formulário de criação no \texttt{SolicitationModal} e reutilizar
o \texttt{PreviewDocumentDialog} sem alterações
(Seção~\ref{subsec:componentes-vue}). A comunicação entre estado e interface ocorre
por diretivas: \texttt{v-model} sustenta os campos do formulário; \texttt{v-if} controla
a exibição condicional do campo ``Nova Categoria''; e \texttt{v-for} percorre a lista de
solicitações na renderização dos \textit{cards}.

Por fim, o projeto adota a \textbf{Options API} em vez da Composition API por
compatibilidade com o restante da aplicação, preservando consistência e reduzindo a
curva de aprendizado para manutenção. Assim, a seção justifica que o front-end foi
integrado de forma incremental, apoiado em reatividade, componentes reutilizáveis,
diretivas e Options API.

\subsection{Nuxt}
\label{subsec:nuxt}

O Nuxt é um \textit{meta-framework} sobre Vue.js que adiciona o \textbf{modo universal}:
a aplicação é renderizada no servidor (\textit{Server-Side Rendering}, SSR) e depois
assume o controle no cliente (\textit{Client-Side Rendering}, CSR) \cite{nuxt2024}. O
SSR melhora a primeira renderização percebida e a indexação; o CSR preserva a fluidez de
uma SPA. Esse modelo é adequado ao Banco de Especialistas, que precisa ser acessível e
indexável, mas também oferece formulários e painéis com interações intensas.

O Nuxt organiza a aplicação por diretórios (\texttt{pages/}, \texttt{components/},
\texttt{layouts/}, \texttt{store/} e \texttt{plugins/}), estrutura adotada no módulo com
\texttt{pages/solicitations/}, \texttt{SolicitationModal.vue},
\texttt{store/specialists} e \texttt{store/appeal}. O \textbf{roteamento por arquivos}
mapeia \texttt{index.vue} para \texttt{/solicitations} e \texttt{\_id.vue} para
\texttt{/solicitations/:id}, refletindo a hierarquia da API. O acesso à API é mediado
pelo \texttt{@nuxtjs/axios}, injetado como \texttt{this.\$axios}, o que centraliza
cabeçalhos, tokens de autenticação e interceptadores de erro.

\subsection{RBAC e Controle de Acesso}
\label{subsec:rbac-controle-acesso}

O controle de acesso baseado em papéis (\textit{Role-Based Access Control}, RBAC) é um
modelo de segurança que restringe o acesso a recursos com base nas funções atribuídas
aos usuários, e não em identidades individuais \cite{sandhu1996}. A ideia central é a
separação entre \textbf{usuários}, \textbf{papéis} e \textbf{permissões}: usuários são
associados a papéis; papéis agregam permissões; e permissões autorizam operações sobre
recursos. Essa indireção simplifica a administração de segurança, porque a concessão ou
revogação de um privilégio é feita sobre o papel, e não sobre cada usuário
individualmente. Em um sistema que lida com dados de pessoas e decisões administrativas,
como o Banco de Especialistas, esse modelo é particularmente adequado: os perfis de
acesso são poucos e bem definidos, e a política de autorização precisa ser auditável.

A escolha do RBAC também se justifica pelo alinhamento com o OWASP Top 10. A edição mais
recente do ranking mantém o \textit{Broken Access Control} na primeira posição dos
riscos mais críticos para aplicações web \cite{owasp2025}, o que reforça a necessidade de
um modelo de autorização explícito, verificado no servidor e auditável. O módulo adota
essa orientação em duas camadas complementares. A camada \textbf{autoritativa} está no
servidor (Seção~\ref{subsec:owasp-seguranca-aplicacoes-web}): cada rota declara o perfil exigido por meio
de um decorador (\texttt{@auth.login\_required}), e o serviço verifica se o solicitante
tem acesso ao recurso específico. A camada de \textbf{interface} apenas reflete, para o
usuário, o que o servidor já garante — ocultando operações para as quais o usuário não
tem permissão, mas sem substituir a verificação no \textit{backend}. Essa divisão é
essencial: a ocultação de um botão é conveniência de usabilidade, não mecanismo de
segurança; um usuário sem perfil que tente chamar a rota diretamente recebe \texttt{403}
do servidor, preservando a integridade do fluxo.

No módulo, os papéis não são mutuamente exclusivos: um mesmo usuário pode ser
solicitante e, ao mesmo tempo, possuir a permissão de revisão, dependendo dos perfis e
das \textit{features} associados à sua conta. A autorização é expressa, portanto, não
apenas por papéis, mas pela combinação entre \textbf{perfil}
(\texttt{ProfileEnum.ADMIN}) e \textbf{permissão}
(\texttt{FeatureEnum.SPECIALIST\_REVIEW}). O Quadro~\ref{quadro:papeis-permissoes}
sintetiza as operações disponíveis no módulo e as condições que cada uma exige.

\begin{quadro}
\centering
\caption{Papéis e permissões no módulo de solicitações.}
\label{quadro:papeis-permissoes}
\begin{tabular}{lcccc}
\hline
\textbf{Operação} & \textbf{Solicitante} & \textbf{Avaliador} & \textbf{Administrador} & \textbf{Requisito} \\
\hline
Criar solicitação & Sim & Sim & Sim & RF01--RF03 \\
Listar próprias solicitações & Sim & Sim & Sim & RF04 \\
Listar todas as solicitações & Não & Sim & Sim & RF04 \\
Consultar detalhe (própria) & Sim & Sim & Sim & RF04 \\
Consultar detalhe (qualquer) & Não & Sim & Sim & RF04 \\
Aprovar ou rejeitar & Não & Sim & Sim & RF05--RF07 \\
Reverter alteração aprovada & Não & Não & Sim & RF09 \\
\hline
\end{tabular}
\fonte{Elaboração própria (2026).}
\end{quadro}

A definição dos papéis operacionais é a seguinte. O \textbf{solicitante} é qualquer
usuário autenticado que cria solicitações e acompanha as suas próprias; não tem acesso à
listagem geral nem à decisão de revisão. O \textbf{avaliador} é o usuário que possui a
\textit{feature} \texttt{SPECIALIST\_REVIEW}, o que lhe permite listar todas as
solicitações, consultar o detalhe de qualquer uma e proferir decisão de aprovação ou
rejeição. O \textbf{administrador} acumula o perfil \texttt{ProfileEnum.ADMIN} e a
\textit{feature} de revisão, e é o único habilitado a executar a reversão de uma
alteração aprovada — operação de efeito sobre o cadastro do especialista que exige, pela
sua gravidade, o nível mais alto de autorização. A tabela também explicita o que cada
papel \textit{não} pode fazer: a coluna de requisito permite rastrear, para cada
operação, o requisito funcional correspondente no
Quadro~\ref{qua:rf} da Seção~\ref{subsec:requisitos-funcionais}, o
que fecha a cadeia de rastreabilidade entre requisito, papel, operaçãoe verificação.
